\documentclass[12pt,a4paper]{article}

% ============================================================
% Dan Koe "How to Fix Your Entire Life in 1 Day" 教學講義
% 作者: 謝慕揚醫師 MD, PhD, FESC
% 來源: X (Twitter) @thedankoe
% 日期: 2026年1月13日
% ============================================================

% Drake_preamble.tex - LaTeX Preamble Template
% 作者: 謝慕揚, MD, PhD, FESC
% 
% 使用說明:
% 1. 表格請使用 longtable，不要有垂直分隔線 |
% 2. 表格內換行使用 \newline，不要使用 \\
% 3. 藥物名稱、檢驗、檢查請維持英文
% 4. Reference 格式請使用 \href 加上驗證過的 DOI

% Including essential packages for Chinese typesetting
\usepackage{xeCJK}
\usepackage{fontspec}
% Configuring Chinese font with Noto Serif CJK TC, fallback to SimSun
\IfFontExistsTF{Noto Serif CJK TC}{
  \setCJKmainfont{Noto Serif CJK TC}
  \setCJKsansfont{Noto Serif CJK TC}
  \setCJKmonofont{Noto Serif CJK TC}
}{\setCJKmainfont{SimSun}
  \setCJKsansfont{SimSun}
  \setCJKmonofont{SimSun}
}

% Including other necessary packages
\usepackage{geometry}
\usepackage{titlesec}
\usepackage{enumitem}
\usepackage{xcolor}
\usepackage{colortbl}
\usepackage{tcolorbox}
\usepackage{booktabs}
\usepackage{longtable}
\usepackage{array}
\usepackage{graphicx}
\usepackage{tikz}
\usepackage{fancyhdr}
\usepackage{multicol}
\usepackage{datetime2}
\usepackage{hyperref}
\usepackage{mdframed}
\usepackage{amsmath}
\usepackage{amssymb}
\usepackage{multirow}

\hypersetup{
    colorlinks=true,
    linkcolor=emergency_blue,
    citecolor=emergency_blue,
    urlcolor=emergency_blue,
    pdfborder={0 0 0}
}

% Defining page geometry
\geometry{
  top=2.5cm,
  bottom=2.5cm,
  left=2cm,
  right=2cm,
  headheight=25pt
}

% Adjusting topmargin to compensate for increased headheight
\addtolength{\topmargin}{-10pt}

% Defining colors
\definecolor{emergency_red}{RGB}{186,24,27}
\definecolor{emergency_blue}{RGB}{0,114,188}
\definecolor{emergency_gray}{RGB}{120,120,120}
\definecolor{highlight_yellow}{RGB}{255,250,205}

% Setting title formats
\titleformat{\section}[block]{\Large\bfseries\color{emergency_red}}{\thesection}{10pt}{}
\titleformat{\subsection}[block]{\large\bfseries\color{emergency_blue}}{\thesubsection}{8pt}{}
\titleformat{\subsubsection}[block]{\normalsize\bfseries\color{emergency_gray}}{\thesubsubsection}{6pt}{}

% Defining custom environment for important notes
\newmdenv[
    linecolor=emergency_red,
    backgroundcolor=highlight_yellow,
    linewidth=2pt,
    topline=true,
    bottomline=true,
    leftline=true,
    rightline=true,
    innertopmargin=10pt,
    innerbottommargin=10pt,
    innerleftmargin=10pt,
    innerrightmargin=10pt
]{importantbox}

% Defining note command with tcolorbox
\newcommand{\note}[1]{
    \par
    \noindent
    \begin{tcolorbox}[
        colback=emergency_blue!5!white,
        colframe=emergency_blue,
        title=\textbf{重點提醒},
        fonttitle=\bfseries,
        boxrule=0.5pt,
        arc=3pt,
        left=6pt,
        right=6pt,
        top=6pt,
        bottom=6pt
    ]
    #1
    \end{tcolorbox}
}

% Key findings box
\newcommand{\keyfinding}[1]{
    \par
    \noindent
    \begin{tcolorbox}[
        colback=yellow!10!white,
        colframe=orange!75!black,
        title=\textbf{核心發現摘要},
        fonttitle=\bfseries,
        boxrule=0.5pt,
        arc=3pt
    ]
    #1
    \end{tcolorbox}
}

% Defining custom date format for ISO 8601
\DTMnewdatestyle{iso}{
  \renewcommand{\DTMdisplaydate}[4]{\number##1-\number##2-\number##3}
  \renewcommand{\DTMDisplaydate}[4]{\number##1-\number##2-\number##3}
}
\DTMsetdatestyle{iso}
\newcommand{\mydate}{\DTMtoday}


% Custom header for this document
\pagestyle{fancy}
\fancyhf{}
\fancyhead[L]{\textcolor{emergency_red}{Dan Koe 個人成長講義}}
\fancyhead[R]{\textcolor{emergency_blue}{行為改變與身份認同}}
\fancyfoot[C]{\textcolor{emergency_gray}{\thepage}}
\renewcommand{\headrulewidth}{1pt}
\renewcommand{\headrule}{\hbox to\headwidth{\color{emergency_red}\leaders\hrule height \headrulewidth\hfill}}

\begin{document}

% ============================================================
% 標題頁 - 社群媒體風格
% ============================================================

\begin{center}
{\LARGE\bfseries\textcolor{emergency_red}{X (Twitter) 長文精選}}\\[0.5cm]
{\Large\textbf{How to Fix Your Entire Life in 1 Day}}\\[0.3cm]
{\large 如何在一天內重新定位你的人生}\\[0.5cm]
{\normalsize 作者：Dan Koe (@thedankoe)}\\[0.2cm]
{\normalsize \href{https://x.com/thedankoe/status/2010751592346030461}{原文連結：X Post 2026年1月13日}}\\[0.3cm]
{\normalsize 整理：謝慕揚 MD, PhD, FESC \quad 日期：\mydate}
\end{center}

\vspace{0.5cm}

% ============================================================
% 核心訊息摘要
% ============================================================

\begin{tcolorbox}[colback=yellow!10!white,colframe=orange!75!black,title=\textbf{核心訊息摘要}]
\textbf{三大核心觀點：}
\begin{enumerate}[leftmargin=*]
    \item \textbf{身份優先改變 (Identity-First Change)}：真正的改變不是改變行為，而是改變「你是誰」
    \item \textbf{目標導向行為 (Teleological Behavior)}：所有行為都是目標導向的，包括你想要逃避的行為
    \item \textbf{智慧即掌舵能力 (Cybernetics)}：智慧是「獲得人生所想」的能力，而非傳統的IQ定義
\end{enumerate}

\textbf{醫療人員應用：}本文對住院醫師職涯發展、避免職業倦怠、建立專業身份認同具有重要啟發。
\end{tcolorbox}

\tableofcontents
\newpage

% ============================================================
\section{前言：為何新年願望總是失敗}
% ============================================================

Dan Koe 開宗明義指出：大多數人的新年目標會失敗，不是因為缺乏紀律，而是因為他們用錯誤的方式嘗試改變。

\subsection{表面目標 vs 深層改變}

設定目標有兩個層次的需求：

\begin{enumerate}
    \item \textbf{改變行為 (Changing Actions)}：次要的、第二順位的改變
    \item \textbf{改變身份 (Changing Who You Are)}：主要的、第一順位的改變
\end{enumerate}

\note{
大多數人設定表面目標，靠意志力維持幾週，然後毫無掙扎地回到舊習慣——因為他們試圖在腐爛的地基上建造美好人生。
}

\subsection{對醫療人員的意義}

對於住院醫師與醫療專業人員而言，這個概念尤其重要：

\begin{itemize}
    \item 「我想要更好的 work-life balance」vs「我是一個懂得設定界限的人」
    \item 「我要多讀 paper」vs「我是一個終身學習者」
    \item 「我不要再 burnout」vs「我是一個懂得照顧自己的醫療專業人員」
\end{itemize}

% ============================================================
\section{第一部分：身份認同與行為改變}
% ============================================================

% ------------------------------------------------------------
\subsection{概念一：你不在想去的地方，是因為你不是會在那裡的人}
% ------------------------------------------------------------

\begin{importantbox}
\textbf{核心洞見：}如果你想要特定的人生結果，你必須在達成目標之前，就先擁有會創造那個結果的生活方式。
\end{importantbox}

想像一個成功的人——可能是體態極佳的健身者、價值數億的創業家、或是能自在社交的魅力人士。

\textbf{關鍵問題：}
\begin{itemize}
    \item 健身者需要「硬撐」才能吃得健康嗎？
    \item CEO 需要強迫自己起床工作嗎？
\end{itemize}

答案是：他們無法想像自己過著其他生活方式。健身者要「硬撐」才能吃不健康；CEO 要「強迫自己」才能賴床——而且他們討厭每一秒。

\note{
當你真正改變了自己，所有不利於目標的習慣都會變得令人厭惡，因為你深刻意識到那些行為會累積成什麼樣的人生。
}

% ------------------------------------------------------------
\subsection{概念二：你不在想去的地方，是因為你不「想」在那裡}
% ------------------------------------------------------------

\begin{quote}
\textit{「只相信行動。生命發生在事件層面，不是言語層面。相信行動。」}\\
\hfill — Alfred Adler
\end{quote}

\subsubsection{目的論行為 (Teleological Behavior)}

理解心智運作的第一步：\textbf{所有行為都是目標導向的}。

\begin{itemize}
    \item 你向前走一步——因為你想到達某個位置
    \item 你抓癢——因為你想消除搔癢感
\end{itemize}

但更深層的是：\textbf{你的目標通常是無意識的}。

\subsubsection{無意識目標的例子}

\begin{longtable}{p{5cm}p{9cm}}
\toprule
\textbf{表面行為} & \textbf{隱藏的無意識目標} \\
\midrule
\endfirsthead

\toprule
\textbf{表面行為} & \textbf{隱藏的無意識目標} \\
\midrule
\endhead

\bottomrule
\endfoot

拖延工作 & 保護自己免於完成作品後被評判的恐懼 \\
\midrule
不離開死氣沉沉的工作 & 追求安全感、可預測性、以及不被視為失敗者的藉口 \\
\midrule
白天躺在沙發上 & 在下一個責任前消磨時間 \\
\midrule
（醫療人員）不設定界限 & 維持「好醫師」的形象、害怕被同事批評 \\
\midrule
（住院醫師）不敢提問 & 保護自己不被認為「不夠聰明」 \\
\bottomrule
\end{longtable}

\begin{tcolorbox}[colback=emergency_blue!5!white, colframe=emergency_blue, title=\textbf{關鍵教訓}]
真正的改變需要改變你的目標——不是設定表面目標，而是改變你的觀點。因為目標就是一種投射到未來的感知濾鏡。
\end{tcolorbox}

% ------------------------------------------------------------
\subsection{概念三：你不在想去的地方，是因為你害怕在那裡}
% ------------------------------------------------------------

\begin{quote}
\textit{「重要的是要記住：無論你是從哪裡得到這個想法，無論你是否見過專業催眠師。但如果你接受了一個想法——來自你自己、老師、父母、朋友、廣告——而且你堅信這個想法是真的，它就對你有催眠師對被催眠者同樣的力量。」}\\
\hfill — Maxwell Maltz
\end{quote}

\subsubsection{身份認同形成的八個步驟}

\begin{longtable}{cp{12cm}}
\toprule
\textbf{步驟} & \textbf{說明} \\
\midrule
\endfirsthead

\multicolumn{2}{c}{\textbf{表格續接}} \\
\toprule
\textbf{步驟} & \textbf{說明} \\
\midrule
\endhead

\bottomrule
\endfoot

1 & \textbf{你想達成一個目標} \\
\midrule
2 & \textbf{你透過該目標的濾鏡感知現實} \\
\midrule
3 & \textbf{你只注意到有助於達成目標的「重要」資訊和想法}（學習） \\
\midrule
4 & \textbf{你朝目標行動並獲得正在進步的回饋} \\
\midrule
5 & \textbf{你重複該行為直到它變成自動化和無意識}（制約） \\
\midrule
6 & \textbf{該行為成為你認為的「自己」的一部分}（「我是那種會...的人」） \\
\midrule
7 & \textbf{你防衛你的身份以維持心理一致性} \\
\midrule
8 & \textbf{你的身份塑造新目標，重啟循環}——如果身份不利於好生活，這會很快惡化 \\
\bottomrule
\end{longtable}

\subsubsection{身份威脅與戰或逃反應}

\begin{itemize}
    \item 當你的身體感到威脅 → 戰或逃反應 (Fight or Flight)
    \item 當你的身份感到威脅 → \textbf{同樣的反應}
\end{itemize}

如果你強烈認同某個政治意識形態，當有人挑戰你的信念時，你會感到壓力——情緒上就像被打了一巴掌。

\note{
對醫療人員而言：如果你無意識地認同「律師」、「遊戲玩家」、或「不會成功的人」，你會避免採取邁向更好生活的行動。同樣地，如果你過度認同「完美醫師」的身份，你可能無法接受自己需要休息或有極限。
}

% ============================================================
\section{第二部分：心智發展框架}
% ============================================================

% ------------------------------------------------------------
\subsection{概念四：你想要的人生存在於特定的心智層次}
% ------------------------------------------------------------

心智會隨時間經歷可預測的發展階段。這已在多個模型中被記錄：Maslow 需求層次、Greuter 自我發展階段、螺旋動力學 (Spiral Dynamics)、整合理論 (Integral Theory)。

\subsubsection{自我發展九階段}

\begin{longtable}{clp{9cm}}
\toprule
\textbf{階段} & \textbf{英文名稱} & \textbf{說明與範例} \\
\midrule
\endfirsthead

\multicolumn{3}{c}{\textbf{表格續接}} \\
\toprule
\textbf{階段} & \textbf{英文名稱} & \textbf{說明與範例} \\
\midrule
\endhead

\bottomrule
\endfoot

1 & Impulsive\newline 衝動期 & 衝動與行動之間沒有分離。黑白思考。\newline \textit{範例：幼兒生氣時打人，因為感覺和行為是同一件事。} \\
\midrule
2 & Self-Protective\newline 自我保護期 & 世界是危險的，學會保護自己。\newline \textit{範例：小孩學會藏成績單、對家事撒謊、說大人想聽的話。} \\
\midrule
3 & Conformist\newline 從眾期 & 你就是你的群體，群體規則感覺像是現實本身。\newline \textit{範例：無法理解為何有人會投給不同政黨。} \\
\midrule
4 & Self-Aware\newline 自我覺察期 & 你注意到內在生活與外在表現不一致。\newline \textit{範例：在教堂裡意識到自己不確定是否相信周圍人相信的，但不知道該怎麼辦。} \\
\midrule
5 & Conscientious\newline 良心期 & 你建立自己的原則系統並對自己負責。\newline \textit{範例：經過仔細研究後離開家庭宗教，採納可以辯護的個人哲學。} \\
\midrule
6 & Individualist\newline 個人主義期 & 你看到自己的原則是被環境塑造的，開始更鬆散地持有它們。\newline \textit{範例：意識到政治觀點更多是因為成長地點而非客觀真理。} \\
\midrule
7 & Strategist\newline 策略期 & 你在系統中工作，同時意識到自己參與其中。\newline \textit{範例：領導組織同時積極質疑自己的盲點。} \\
\midrule
8 & Construct-Aware\newline 建構覺察期 & 你看到所有框架，包括身份，都是有用的虛構。\newline \textit{範例：以隱喻而非字面意義持有精神信仰，知道地圖不是領土。} \\
\midrule
9 & Unitive\newline 統一期 & 自我與生命之間的分離消融。\newline \textit{範例：工作、休息、玩樂感覺是同一件事。沒有人需要成為什麼。} \\
\bottomrule
\end{longtable}

\note{
大多數讀者可能在第4到第8階段之間。好消息是：無論你在哪個階段，向前移動都遵循相同模式。
}

% ------------------------------------------------------------
\subsection{概念五：智慧是獲得人生所想的能力}
% ------------------------------------------------------------

\begin{quote}
\textit{「智慧的唯一真正測試是：你是否得到了你想要的人生。」}\\
\hfill — Naval Ravikant
\end{quote}

\subsubsection{成功公式}

\begin{center}
\fbox{\Large \textbf{成功 = 主動性 (Agency) + 機會 (Opportunity) + 智慧 (Intelligence)}}
\end{center}

\begin{itemize}
    \item 高主動性但低機會 → 無論你多麼可能採取行動，目標都不會有太多成果
    \item 有機會和主動性但低智慧 → 永遠無法完全受益於機會
\end{itemize}

\subsubsection{控制論 (Cybernetics)}

Cybernetics 來自希臘字 \textit{kybernetikos}，意為「掌舵」或「善於掌舵」。

\textbf{智慧系統的特性：}
\begin{enumerate}
    \item 有一個目標
    \item 朝目標行動
    \item 感知自己在哪裡
    \item 與目標比較
    \item 根據回饋再次行動
\end{enumerate}

\begin{tcolorbox}[colback=emergency_blue!5!white, colframe=emergency_blue, title=\textbf{高智慧 vs 低智慧}]
\textbf{高智慧：}迭代、堅持、理解大局的能力。任何問題在足夠長的時間尺度上都可以解決。

\textbf{低智慧的標誌：}無法從錯誤中學習。遇到路障就放棄。
\end{tcolorbox}

\subsubsection{如何變得更有智慧}

\begin{enumerate}
    \item 拒絕已知路徑
    \item 潛入未知
    \item 設定更高的新目標以擴展心智
    \item 擁抱混沌並允許成長
    \item 研究自然的普遍原則
    \item 成為深度通才 (Deep Generalist)
\end{enumerate}

% ============================================================
\section{第三部分：一日重啟協議}
% ============================================================

\begin{importantbox}
\textbf{協議目標：}幫助你達到不協調點、穿越不確定性、發現你真正想要追求的事物——清晰到壓倒性，讓干擾不再有分量。
\end{importantbox}

Dan Koe 觀察到成功翻轉身份的人會經歷三個階段：

\begin{enumerate}
    \item \textbf{不協調 (Dissonance)}：感覺不屬於當前生活，對缺乏進步感到厭倦
    \item \textbf{不確定 (Uncertainty)}：不知道接下來是什麼，要嘛實驗要嘛迷失
    \item \textbf{發現 (Discovery)}：發現想要追求的事物，在6個月內取得6年的進步
\end{enumerate}

% ------------------------------------------------------------
\subsection{Part 1 早晨：心理挖掘 — 願景與反願景}
% ------------------------------------------------------------

設定 15-30 分鐘思考並回答以下問題。\textbf{不要將這個思考外包給 AI。}

\subsubsection{覺察當前痛苦的問題}

\begin{enumerate}
    \item 你學會忍受的沉悶持續不滿是什麼？（不是深層痛苦，而是你學會容忍的）
    \item 你反覆抱怨但從未真正改變的事情是什麼？寫下過去一年你最常表達的三個抱怨。
    \item 對於每個抱怨：如果有人觀察你的行為（不是言語），他們會認為你實際上想要什麼？
    \item 關於你當前生活，什麼真相是你無法承認給你深深尊敬的人聽的？
\end{enumerate}

\subsubsection{建立反願景 (Anti-Vision)}

\begin{enumerate}
    \item 如果未來五年什麼都不變，描述一個普通的週二。你在哪裡醒來？身體感覺如何？第一個想到什麼？誰在你身邊？早上9點到下午6點做什麼？晚上10點感覺如何？
    \item 現在對十年做同樣的事。你錯過了什麼？哪些機會關閉了？誰放棄了你？當你不在房間時人們怎麼說你？
    \item 你在生命盡頭。你過了安全版本的人生。你從未打破模式。代價是什麼？你從未讓自己感受、嘗試或成為什麼？
    \item 你生活中誰已經在過你剛描述的未來？比你早5、10、20年在同一軌道上的人？當你想到變成他們時感覺如何？
    \item 你必須放棄什麼身份才能真正改變？（「我是那種...的人」）不再是那個人會讓你付出什麼社交代價？
    \item 你沒有改變的最尷尬原因是什麼？讓你聽起來軟弱、害怕或懶惰而非合理的原因？
    \item 如果你當前的行為是一種自我保護，你到底在保護什麼？這種保護讓你付出什麼代價？
\end{enumerate}

\subsubsection{建立願景 (Vision)}

\begin{enumerate}
    \item 暫時忘記實際性。如果你可以彈指間在三年後過著不同的生活——不是什麼是現實的，而是你真正想要什麼？描述一個普通的週二，和問題5同樣的細節程度。
    \item 你必須對自己相信什麼，那種生活才會感覺自然而非強迫？寫下身份宣言：「我是那種...的人」
    \item 如果你已經是那個人，這週你會做什麼事？
\end{enumerate}

% ------------------------------------------------------------
\subsection{Part 2 白天：打斷自動駕駛模式}
% ------------------------------------------------------------

全天不斷思考你在早晨日誌中寫的內容。設定手機提醒或日曆事件，包含問題以便立即開始思考。

\textbf{建議提醒時間：}
\begin{itemize}
    \item \textbf{11:00am}：我現在做的事情是在逃避什麼？
    \item \textbf{1:30pm}：如果有人拍攝過去兩小時，他們會認為我想要什麼樣的人生？
    \item \textbf{3:15pm}：我正在走向我討厭的生活還是想要的生活？
    \item \textbf{5:00pm}：我假裝不重要的最重要的事是什麼？
    \item \textbf{7:30pm}：今天我做了什麼是出於身份保護而非真正渴望？（提示：大多數事情）
    \item \textbf{9:00pm}：今天我什麼時候感覺最有活力？什麼時候感覺最死氣沉沉？
\end{itemize}

\textbf{額外思考問題（通勤、散步或躺著時）：}
\begin{itemize}
    \item 如果我不再需要人們把我看作 [你在問題10寫的身份]，會改變什麼？
    \item 我生活中哪裡在用安全交換活力？
    \item 我想成為的人的最小版本是什麼，我明天就可以做到？
\end{itemize}

% ------------------------------------------------------------
\subsection{Part 3 晚間：整合洞見}
% ------------------------------------------------------------

\begin{enumerate}
    \item 經過今天，關於你為何停滯不前，什麼感覺最真實？
    \item 真正的敵人是什麼？清楚地說出來。不是環境。不是其他人。一直在主導的內在模式或信念。
    \item 寫一句話捕捉你拒絕讓生活變成什麼。這是你的反願景壓縮版。讀它時應該讓你有感覺。
    \item 寫一句話捕捉你正在建立什麼，知道它會演變。這是你的願景MVP。
\end{enumerate}

\subsubsection{設定目標透鏡}

\begin{enumerate}
    \item \textbf{一年透鏡：}一年後什麼必須成真，你才會知道你打破了舊模式？一件具體的事。
    \item \textbf{一月透鏡：}一個月後什麼必須成真，一年透鏡才可能保持？
    \item \textbf{每日透鏡：}明天你可以時間區塊的2-3個行動是什麼，是你正在成為的那個人會簡單地做的？
\end{enumerate}

% ============================================================
\section{【實作練習】反思問題工作表}
% ============================================================

\textit{請用筆填寫以下工作表。這個練習需要約一整天完成，但效果將持續更久。}

% ------------------------------------------------------------
\subsection{早晨：心理挖掘區}
% ------------------------------------------------------------

\textbf{1. 我學會容忍的沉悶不滿是什麼？}
\begin{center}
\fbox{\parbox{0.9\textwidth}{\vspace{3cm}}}
\end{center}

\textbf{2. 過去一年我最常抱怨的三件事：}
\begin{center}
\fbox{\parbox{0.9\textwidth}{
\begin{enumerate}
    \item \vspace{1.5cm}
    \item \vspace{1.5cm}
    \item \vspace{1.5cm}
\end{enumerate}
}}
\end{center}

\textbf{3. 觀察我的行為，別人會認為我實際上想要什麼？}
\begin{center}
\fbox{\parbox{0.9\textwidth}{\vspace{3cm}}}
\end{center}

\textbf{4. 我無法承認給尊敬的人聽的真相是什麼？}
\begin{center}
\fbox{\parbox{0.9\textwidth}{\vspace{3cm}}}
\end{center}

\newpage

\textbf{5. 如果未來五年什麼都不變，我普通的週二會是什麼樣子？}
\begin{center}
\fbox{\parbox{0.9\textwidth}{\vspace{5cm}}}
\end{center}

\textbf{6. 我必須放棄什麼身份才能真正改變？}
\begin{center}
\fbox{\parbox{0.9\textwidth}{
「我是那種 \underline{\hspace{8cm}} 的人」\\[1cm]
社交代價：\vspace{2cm}
}}
\end{center}

\textbf{7. 我沒有改變的最尷尬原因（聽起來軟弱/害怕/懶惰的）：}
\begin{center}
\fbox{\parbox{0.9\textwidth}{\vspace{3cm}}}
\end{center}

\textbf{8. 三年後我真正想要的生活（普通的週二）：}
\begin{center}
\fbox{\parbox{0.9\textwidth}{\vspace{5cm}}}
\end{center}

\textbf{9. 我的新身份宣言：}
\begin{center}
\fbox{\parbox{0.9\textwidth}{
\LARGE 「我是那種 \underline{\hspace{6cm}} 的人」
\vspace{1cm}
}}
\end{center}

\newpage

% ------------------------------------------------------------
\subsection{白天：自動駕駛中斷提醒設定}
% ------------------------------------------------------------

請現在在手機上設定以下提醒：

\begin{longtable}{p{3cm}p{11cm}}
\toprule
\textbf{時間} & \textbf{提醒內容} \\
\midrule
\endfirsthead

\bottomrule
\endfoot

11:00am & 我現在做的事情是在逃避什麼？\\
& \fbox{\parbox{9cm}{\vspace{1.5cm}}} \\
\midrule
1:30pm & 如果有人拍攝過去兩小時，他們會認為我想要什麼？\\
& \fbox{\parbox{9cm}{\vspace{1.5cm}}} \\
\midrule
3:15pm & 我正在走向我討厭的生活還是想要的生活？\\
& \fbox{\parbox{9cm}{\vspace{1.5cm}}} \\
\midrule
5:00pm & 我假裝不重要的最重要的事是什麼？\\
& \fbox{\parbox{9cm}{\vspace{1.5cm}}} \\
\midrule
7:30pm & 今天我做了什麼是出於身份保護而非真正渴望？\\
& \fbox{\parbox{9cm}{\vspace{1.5cm}}} \\
\midrule
9:00pm & 今天我什麼時候感覺最有活力？最死氣沉沉？\\
& \fbox{\parbox{9cm}{\vspace{1.5cm}}} \\
\bottomrule
\end{longtable}

% ------------------------------------------------------------
\subsection{晚間：整合與行動規劃}
% ------------------------------------------------------------

\textbf{1. 經過今天，我停滯不前的真正原因是：}
\begin{center}
\fbox{\parbox{0.9\textwidth}{\vspace{3cm}}}
\end{center}

\textbf{2. 真正的敵人（內在模式或信念）：}
\begin{center}
\fbox{\parbox{0.9\textwidth}{\vspace{2cm}}}
\end{center}

\textbf{3. 我的反願景壓縮版（一句話）：}
\begin{center}
\fbox{\parbox{0.9\textwidth}{
我拒絕讓我的生活變成：\vspace{1.5cm}
}}
\end{center}

\textbf{4. 我的願景MVP（一句話）：}
\begin{center}
\fbox{\parbox{0.9\textwidth}{
我正在建立：\vspace{1.5cm}
}}
\end{center}

\newpage

% ------------------------------------------------------------
\subsection{六大要素規劃表}
% ------------------------------------------------------------

\begin{longtable}{p{4cm}p{10cm}}
\toprule
\textbf{要素} & \textbf{你的答案} \\
\midrule
\endfirsthead

\multicolumn{2}{c}{\textbf{表格續接}} \\
\toprule
\textbf{要素} & \textbf{你的答案} \\
\midrule
\endhead

\bottomrule
\endfoot

\textbf{Anti-Vision}\newline 反願景\newline {\small 我絕不想再經歷的生活} & \fbox{\parbox{8.5cm}{\vspace{3cm}}} \\
\midrule
\textbf{Vision}\newline 願景\newline {\small 我想要並可改進的理想生活} & \fbox{\parbox{8.5cm}{\vspace{3cm}}} \\
\midrule
\textbf{1 Year Goal}\newline 一年目標\newline {\small 一年後打破舊模式的具體證據} & \fbox{\parbox{8.5cm}{\vspace{2.5cm}}} \\
\midrule
\textbf{1 Month Project}\newline 一月專案\newline {\small 需要學習什麼？建立什麼？} & \fbox{\parbox{8.5cm}{\vspace{2.5cm}}} \\
\midrule
\textbf{Daily Levers}\newline 每日槓桿\newline {\small 推動專案完成的優先任務} & \fbox{\parbox{8.5cm}{\vspace{2.5cm}}} \\
\midrule
\textbf{Constraints}\newline 限制條件\newline {\small 我不願意犧牲的是什麼} & \fbox{\parbox{8.5cm}{\vspace{2.5cm}}} \\
\bottomrule
\end{longtable}

% ============================================================
\section{第四部分：遊戲化人生系統}
% ============================================================

\begin{quote}
\textit{「內在體驗的最佳狀態是意識中有秩序。當心理能量——或注意力——投資於現實目標，且技能與行動機會匹配時，這就會發生。追求目標在意識中帶來秩序，因為人必須專注於手頭的任務，暫時忘記其他一切。」}\\
\hfill — Mihaly Csikszentmihalyi
\end{quote}

\subsection{為什麼遊戲讓人上癮？}

遊戲是痴迷、享受和心流狀態 (Flow State) 的典範。它們擁有帶來專注和清晰的所有組件。如果我們逆向工程這些組件，就可以活在更深享受、更少干擾、更多成功的狀態中。

\subsection{遊戲隱喻對照表}

\begin{longtable}{p{4cm}p{4cm}p{6cm}}
\toprule
\textbf{遊戲元素} & \textbf{人生對應} & \textbf{功能} \\
\midrule
\endfirsthead

\toprule
\textbf{遊戲元素} & \textbf{人生對應} & \textbf{功能} \\
\midrule
\endhead

\bottomrule
\endfoot

\textbf{How you win}\newline 如何獲勝 & Vision 願景 & 至少在遊戲演變前，這是你的勝利定義 \\
\midrule
\textbf{What's at stake}\newline 賭注是什麼 & Anti-Vision 反願景 & 如果你輸了或放棄會發生什麼 \\
\midrule
\textbf{The mission}\newline 任務 & 1 Year Goal\newline 一年目標 & 你人生中唯一的優先事項 \\
\midrule
\textbf{The boss fight}\newline Boss 戰 & 1 Month Project\newline 一月專案 & 你如何獲得經驗值和裝備 \\
\midrule
\textbf{The quests}\newline 任務線 & Daily Levers\newline 每日槓桿 & 解鎖新機會的日常流程 \\
\midrule
\textbf{The rules}\newline 規則 & Constraints\newline 限制條件 & 鼓勵創造力的限制 \\
\bottomrule
\end{longtable}

\begin{tcolorbox}[colback=emergency_blue!5!white, colframe=emergency_blue, title=\textbf{專注力護盾}]
所有這些元素形成一組同心圓，像力場一樣保護你的心智免受干擾和閃亮物體的誘惑。你越玩這個遊戲，這個力量就越強，很快它就會成為你是誰的一部分——你不會想要其他方式。
\end{tcolorbox}

% ============================================================
\section{醫療人員應用建議}
% ============================================================

\subsection{住院醫師職涯發展應用}

\begin{enumerate}
    \item \textbf{身份先於技能}
    \begin{itemize}
        \item 不要只想「我要學會做某個 procedure」
        \item 而是「我是一個持續精進技術的介入心臟科醫師」
    \end{itemize}

    \item \textbf{識別無意識目標}
    \begin{itemize}
        \item 為什麼你避免在晨會提問？（保護自己不被認為無知）
        \item 為什麼你不敢拒絕額外工作？（維持「好 R」的形象）
    \end{itemize}

    \item \textbf{設定一年透鏡}
    \begin{itemize}
        \item 一年後什麼必須成真，我才知道我在正確的軌道上？
        \item 例：「我能獨立處理 Complex PCI」「我有完整的研究計畫」
    \end{itemize}
\end{enumerate}

\subsection{醫療工作中的身份認同}

許多醫師過度認同「完美醫師」「永遠在線」「犧牲奉獻」的身份，這可能導致：
\begin{itemize}
    \item 無法接受自己需要休息
    \item 當病人 outcome 不好時過度自責
    \item 忽視自身健康
\end{itemize}

\note{
健康的醫師身份應該是：「我是一個照顧病人同時也照顧自己的醫療專業人員。」這不是自私——只有健康的醫師才能長期提供高品質照護。
}

\subsection{避免職業倦怠的心理框架}

\begin{enumerate}
    \item \textbf{建立反願景}：想像10年後完全燒盡、對醫學失去熱情的自己
    \item \textbf{建立願景}：想像10年後在專業上有成就、生活有品質的自己
    \item \textbf{設定限制條件}：我不願意犧牲什麼？（家庭時間？運動？興趣？）
    \item \textbf{每日槓桿}：什麼日常行動可以同時推進職業發展和生活品質？
\end{enumerate}

\begin{tcolorbox}[colback=yellow!10!white,colframe=orange!75!black,title=\textbf{給住院醫師的提醒}]
訓練期間會過去。你現在建立的身份認同、習慣和價值觀，會決定你成為什麼樣的主治醫師。不要等到「以後有時間」才開始照顧自己——那個「以後」永遠不會來。
\end{tcolorbox}

% ============================================================
\section{結論與行動呼籲}
% ============================================================

Dan Koe 的這篇文章核心訊息是：

\begin{enumerate}
    \item \textcolor{red}{\textbf{改變從身份開始}}：不是改變你「做」什麼，而是改變你「是」誰
    \item \textcolor{red}{\textbf{所有行為都有目的}}：你拖延、逃避的行為背後都有無意識目標
    \item \textcolor{red}{\textbf{智慧是可以培養的}}：通過迭代、學習、擁抱不確定性來提升
    \item \textcolor{red}{\textbf{遊戲化你的人生}}：建立願景、反願景、目標、專案、每日行動、限制條件的系統
\end{enumerate}

\vspace{0.5cm}

\begin{center}
\fbox{\parbox{0.8\textwidth}{
\centering
\textbf{行動呼籲}\\[0.3cm]
選一天，花15-30分鐘完成本講義的工作表。\\
設定手機提醒，全天打斷自動駕駛模式。\\
晚上整合洞見，開始你的新遊戲。
}}
\end{center}

% ============================================================
\section{參考文獻與延伸閱讀}
% ============================================================

\begin{enumerate}
    \item Koe D. How to fix your entire life in 1 day. \href{https://x.com/thedankoe/status/2010751592346030461}{X (Twitter). 2026年1月13日.}

    \item Adler A. \textit{Understanding Human Nature}. 1927.（阿德勒心理學：目的論行為）

    \item Maltz M. \textit{Psycho-Cybernetics}. 1960.（心理控制論）

    \item Csikszentmihalyi M. \textit{Flow: The Psychology of Optimal Experience}. 1990.（心流理論）

    \item Cook-Greuter S. Ego Development: Nine Levels of Increasing Embrace. 2013.（自我發展九階段）

    \item Maslow AH. \textit{A Theory of Human Motivation}. Psychological Review. 1943.（需求層次理論）

    \item Naval Ravikant. \href{https://nav.al/}{The Almanack of Naval Ravikant.}（智慧與財富哲學）
\end{enumerate}

% ============================================================
\section{術語對照表}
% ============================================================

\begin{longtable}{p{5cm}p{9cm}}
\toprule
\textbf{英文術語} & \textbf{中文翻譯與說明} \\
\midrule
\endfirsthead

\multicolumn{2}{c}{\textbf{表格續接}} \\
\toprule
\textbf{英文術語} & \textbf{中文翻譯與說明} \\
\midrule
\endhead

\bottomrule
\endfoot

Identity & 身份認同：你認為自己是誰 \\
Teleological & 目的論的：所有行為都是目標導向的 \\
Cybernetics & 控制論：研究系統如何自我調節以達成目標 \\
Agency & 主動性：採取行動的能力和意願 \\
Anti-Vision & 反願景：你絕不想要的生活樣貌 \\
Vision & 願景：你想要建立的理想生活 \\
Flow State & 心流狀態：完全沉浸於活動中的最佳體驗狀態 \\
Ego Development & 自我發展：心智經歷的可預測發展階段 \\
Conditioning & 制約：行為因重複而變成自動化的過程 \\
Psychological Consistency & 心理一致性：維持自我認知穩定的傾向 \\
Deep Generalist & 深度通才：在多領域有深度理解的人 \\
Autopilot & 自動駕駛：無意識地按習慣行動 \\
\bottomrule
\end{longtable}

\end{document}
